% v13. Complete fixed-denominator divisor calculation; no denominator cutoff.
\section{小さい分母を持つ二価 branch の局所伝播}
\label{sp13:section}

二高分母の局所模型を共有高 branch の両端で比較するとき，
分母だけでなく係数 $D$ の値を保持することが有効である．
ここでは小さい四つの分母について，可能な係数を上限なしに決定する．
結果は閉路に限らず，より高い次数の vertices に接続する二価部分にも適用できる．

\begin{proposition}[固定した小分母の係数集合]
\label{sp13:coefficient-sets}
奇数の二高分母を持つ正の整数局所模型において，
分母 $n$ の高 ray が取り得る係数 $D$ の集合を $\mathcal D(n)$ と書く．
このとき
\[
 \mathcal D(3)=\{3,6\},\qquad
 \mathcal D(5)=\{6,11,18,21\},
\]
\[
 \mathcal D(7)=F_7\cup\{7k-9:k\geq1\},\qquad
 F_7=\{-7,0,3,7,9,16,21,23,44\},
\]
\[
 \mathcal D(13)=F_{13}\cup\{13k-14:k\geq1\},
\]
ただし
\begin{align*}
 F_{13}={}&\{-15,-13,-7,-6,-3,-2,0,13,14,15,18,22,23,26,\notag\\
          &\hspace{13mm}27,28,31,41,44,52,57,58,80,83,92,161\}.
\end{align*}
特に
\[
 \mathcal D(7)\cap(-\mathcal D(7))=\{-7,0,7\},\qquad
 \mathcal D(13)\cap(-\mathcal D(13))=\{-15,-13,0,13,15\}.
\]
これらは相手の分母に上限を置かない等式である．
\end{proposition}

\begin{proof}
定理~\ref{fp12:partner-structure} の二方向の factorization を用いる．
固定した分母が canonical な $A$ 側の場合には，
$r,s\mid n-1$ と $z\mid|(r+1)(s-(n-1)/r-1)|$ を調べる．
同定理で証明した通り，最後の絶対値の中が零となる枝に解はない．
$B$ 側の場合には，$r+1\mid n-1$ と $z\mid n-1$ を選び，
$C_B=r((n-1)/(r+1)-z-1)\ne0$ なら $s\mid|C_B|$ を調べる．
各候補に正整数性，残る整除条件，奇数性を課す．

非共鳴の固定 ray 付き模型の数は $n=3,5,7,13$ の順に $2,4,13,28$ である．
分母三と五のデータは命題~\ref{fp12:small-partners} に表示されている．
各模型で
\[
 D_A=-a+B-r+S,\qquad D_B=-b+A+r+1-T
\]
を評価すると，それぞれ表示した有限集合を得る．
これは証明された有限個の約数候補を全て評価する計算であり，
相手分母や numerator の打切りを含まない．
全ての候補，棄却理由，受理された標準形と係数は付属の
\texttt{small\_denominator\_propagation\_verification\_v13.json}
で再生成できる．

残る共鳴は既知の二族だけである．
$(a,b,r,s,z,S,T)=(6k+1,7,2,2k,1,3,6)$ では固定 ray は $B$ であり，
$D_B=7k-9$ となる．
$(a,b,r,s,z,S,T)=(12k+1,13,3,3k,2,4,6)$ では $D_B=13k-14$ である．
これで全ての模型が尽くされる．
両集合とその符号反転との共通部分を取ると最後の二式を得る．
\end{proof}

\begin{theorem}[小分母の二価接続に必要な隣接分母]
\label{sp13:propagation}
高 branch の両端の core が二高分母型であり，
正の整数局所模型と対向係数の相殺 $D+D'=0$ を仮定する．
共有分母を $n$ とし，両端の core におけるもう一つの高分母を $u,v$ とする．
このとき次が成り立つ．
\begin{enumerate}
\item $n=3$ または $5$ は不可能である．
\item $n=7$ で $D\ne0$ なら，
      $\{u,v\}=\{3,13\}$ または $\{3,19\}$ であり，係数は $\pm7$ である．
      $D=0$ なら $u=v=7$ である．
\item $n=13$ で $D\ne0$ なら，
      $\{u,v\}=\{5,25\}$ で係数 $\pm15$，または
      $\{u,v\}=\{7,21\}$ で係数 $\pm13$ となる．
      $D=0$ なら $u,v\in\{13,17\}$ である．
\end{enumerate}
ここで零係数の可能性は，閉路などで既に非零性が証明されている場合を除き，
取り除いてはならない．この定理は接着の必要条件であり，
列挙された分母の組に対する global atlas の存在を主張しない．
\end{theorem}

\begin{proof}
命題~\ref{sp13:coefficient-sets} により，三と五では係数が常に正なので相殺できない．
七と十三では同命題の共通部分を用い，残る係数を持つ模型だけを取り出す．
その全データは次の通りである．表の ``側'' は固定分母の canonical な位置を表す．
\begin{center}
\small
\begin{tabular}{ccrrrrrr}
 $n$ & 側 & $a$ & $b$ & $r$ & $s$ & $z$ & $D_n$\\\hline
 7&B&3&7&1&1&1&$-7$\\
 7&A&7&13&2&1&1&7\\
 7&A&7&19&1&2&2&7\\
 7&A&7&7&2&2&1&0\\\hline
 13&B&5&13&1&1&1&$-15$\\
 13&A&13&25&2&1&1&15\\
 13&B&7&13&2&1&1&$-13$\\
 13&A&13&21&3&1&1&13\\
 13&A&13&13&3&3&2&0\\
 13&A&13&17&1&3&1&0
\end{tabular}
\end{center}
共鳴の係数列に新たな相殺可能値はない．
表の同じ $n$ に対して正負が相殺する行を組にすれば結論となる．
\end{proof}

\begin{corollary}[固定七・十三の共鳴枝は二価接続を延長できない]
\label{sp13:no-resonant-continuation}
定理~\ref{fp12:partner-structure} の共鳴模型の固定 ray $b=7$ または $13$ を，
もう一つの二高分母 core に接続し，対向 $D$ を相殺することはできない．
この結論に node 非負性も閉路の仮定も必要ない．
\end{corollary}

\begin{proof}
固定 ray の係数はそれぞれ $7k-9$，$13k-14$，$k\geq1$ である．
これらは命題~\ref{sp13:coefficient-sets} の各共通部分に属さない．
\end{proof}
